% Options for packages loaded elsewhere
\PassOptionsToPackage{unicode}{hyperref}
\PassOptionsToPackage{hyphens}{url}
\documentclass[
  11pt,
  a4paper,
]{article}
\usepackage{xcolor}
\usepackage[margin=18mm]{geometry}
\usepackage{amsmath,amssymb}
\setcounter{secnumdepth}{-\maxdimen} % remove section numbering
\usepackage{iftex}
\ifPDFTeX
  \usepackage[T1]{fontenc}
  \usepackage[utf8]{inputenc}
  \usepackage{textcomp} % provide euro and other symbols
\else % if luatex or xetex
  \usepackage{unicode-math} % this also loads fontspec
  \defaultfontfeatures{Scale=MatchLowercase}
  \defaultfontfeatures[\rmfamily]{Ligatures=TeX,Scale=1}
\fi
\ifPDFTeX\else
  % xetex/luatex font selection
    \setmainfont[]{Noto Serif}
    \setsansfont[]{Noto Sans}
    \setmonofont[]{Noto Sans Mono}
  \setmathfont[]{Noto Sans Math}
\fi
% Use upquote if available, for straight quotes in verbatim environments
\IfFileExists{upquote.sty}{\usepackage{upquote}}{}
\IfFileExists{microtype.sty}{% use microtype if available
  \usepackage[]{microtype}
  \UseMicrotypeSet[protrusion]{basicmath} % disable protrusion for tt fonts
}{}
\makeatletter
\@ifundefined{KOMAClassName}{% if non-KOMA class
  \IfFileExists{parskip.sty}{%
    \usepackage{parskip}
  }{% else
    \setlength{\parindent}{0pt}
    \setlength{\parskip}{6pt plus 2pt minus 1pt}}
}{% if KOMA class
  \KOMAoptions{parskip=half}}
\makeatother
\ifLuaTeX
\usepackage[bidi=basic]{babel}
\else
\usepackage[bidi=default]{babel}
\fi
\babelprovide[main,import]{russian}
\ifPDFTeX
\else
\babelfont{rm}[]{Noto Serif}
\fi
% get rid of language-specific shorthands (see #6817):
\let\LanguageShortHands\languageshorthands
\def\languageshorthands#1{}
\setlength{\emergencystretch}{3em} % prevent overfull lines
\providecommand{\tightlist}{%
  \setlength{\itemsep}{0pt}\setlength{\parskip}{0pt}}
\usepackage{bookmark}
\IfFileExists{xurl.sty}{\usepackage{xurl}}{} % add URL line breaks if available
\urlstyle{same}
\hypersetup{
  pdftitle={Билеты по дискретным случайным процессам},
  pdflang={ru-RU},
  hidelinks,
  pdfcreator={LaTeX via pandoc}}

\title{Билеты по дискретным случайным процессам}
\author{}
\date{}

\begin{document}
\maketitle

\section{Билет 15. Теорема Колмогорова о согласованных конечномерных
распределениях}\label{ux431ux438ux43bux435ux442-15.-ux442ux435ux43eux440ux435ux43cux430-ux43aux43eux43bux43cux43eux433ux43eux440ux43eux432ux430-ux43e-ux441ux43eux433ux43bux430ux441ux43eux432ux430ux43dux43dux44bux445-ux43aux43eux43dux435ux447ux43dux43eux43cux435ux440ux43dux44bux445-ux440ux430ux441ux43fux440ux435ux434ux435ux43bux435ux43dux438ux44fux445}

Источники: Ширяев 1, §9; Сердобольская, лекция 1.

\subsection{1. Короткий
ответ}\label{ux43aux43eux440ux43eux442ux43aux438ux439-ux43eux442ux432ux435ux442}

Теорема Колмогорова отвечает на главный вопрос теории процессов: когда
набор конечномерных распределений действительно является набором
распределений некоторого случайного процесса.

Пусть \(T\) - множество моментов времени, \(E\) - пространство
состояний, например \(E=\mathbb R\), и для каждого конечного набора

\[
t_1,\ldots,t_n\in T
\]

задана вероятностная мера

\[
\mu_{t_1,\ldots,t_n}
\]

на \(E^n\). Эти меры должны быть согласованы:

\begin{enumerate}
\def\labelenumi{\arabic{enumi}.}
\tightlist
\item
  перестановка моментов времени переставляет координаты распределения;
\item
  удаление лишних координат дает соответствующее маргинальное
  распределение.
\end{enumerate}

Тогда существует вероятностная мера \(P\) на пространстве траекторий
\(E^T\) такая, что для координатного процесса

\[
X_t(x)=x(t),\qquad x\in E^T,
\]

его конечномерные распределения равны заданным:

\[
P\{(X_{t_1},\ldots,X_{t_n})\in B\}
=
\mu_{t_1,\ldots,t_n}(B).
\]

Иначе говоря, согласованное семейство конечномерных распределений задает
закон процесса на пространстве траекторий.

\subsection{2. Полный
ответ}\label{ux43fux43eux43bux43dux44bux439-ux43eux442ux432ux435ux442}

\subsubsection{2.1. Введение и
мотивация}\label{ux432ux432ux435ux434ux435ux43dux438ux435-ux438-ux43cux43eux442ux438ux432ux430ux446ux438ux44f}

В Билете 06 конечномерные распределения процесса были введены как
распределения случайных векторов

\[
(X_{t_1},\ldots,X_{t_n}).
\]

Если процесс уже задан, то его конечномерные распределения автоматически
существуют и согласованы между собой. В билете 15 вопрос обратный: можно
ли сначала задать все конечномерные распределения, а потом получить
случайный процесс, для которого они действительно являются
распределениями конечных наборов координат.

Именно это делает теорема Колмогорова о согласованных конечномерных
распределениях. Она является базовой конструкцией существования
случайных процессов. Без нее определения вроде ``пусть есть процесс с
такими конечномерными распределениями'' были бы неполными: нужно
доказать, что такой процесс вообще существует.

\subsubsection{2.2. Основные обозначения: время и фазовое
пространство}\label{ux43eux441ux43dux43eux432ux43dux44bux435-ux43eux431ux43eux437ux43dux430ux447ux435ux43dux438ux44f-ux432ux440ux435ux43cux44f-ux438-ux444ux430ux437ux43eux432ux43eux435-ux43fux440ux43eux441ux442ux440ux430ux43dux441ux442ux432ux43e}

Пусть \(T\) - множество индексов. В дискретном времени обычно

\[
T=\mathbb N
\]

или

\[
T=\mathbb Z.
\]

В непрерывном времени может быть

\[
T=\mathbb R_+
\]

или

\[
T=[0,\infty).
\]

Пусть \(E\) - пространство значений процесса. Для экзамена достаточно
уверенно рассматривать случай

\[
E=\mathbb R
\]

с борелевской \(\sigma\)-алгеброй \(\mathcal B(\mathbb R)\), но
формулировка работает и для стандартных борелевских пространств.

\subsubsection{2.3. Пространство траекторий и цилиндрические
множества}\label{ux43fux440ux43eux441ux442ux440ux430ux43dux441ux442ux432ux43e-ux442ux440ux430ux435ux43aux442ux43eux440ux438ux439-ux438-ux446ux438ux43bux438ux43dux434ux440ux438ux447ux435ux441ux43aux438ux435-ux43cux43dux43eux436ux435ux441ux442ux432ux430}

Пространство всех траекторий:

\[
E^T=\{x:T\to E\}.
\]

Для каждого \(t\in T\) есть координатное отображение

\[
X_t:E^T\to E,
\qquad
X_t(x)=x(t).
\]

Если выбрать конечный набор индексов

\[
t_1,\ldots,t_n,
\]

то есть конечномерная проекция

\[
\pi_{t_1,\ldots,t_n}:E^T\to E^n,
\qquad
\pi_{t_1,\ldots,t_n}(x)=(x(t_1),\ldots,x(t_n)).
\]

Цилиндрическим множеством называется прообраз борелевского множества при
такой конечномерной проекции:

\[
C=\pi_{t_1,\ldots,t_n}^{-1}(B),
\qquad
B\in\mathcal B(E^n).
\]

То есть

\[
C=\{x\in E^T:(x(t_1),\ldots,x(t_n))\in B\}.
\]

Интуитивно цилиндрическое событие смотрит только на конечное число
координат траектории и не ограничивает остальные координаты.

\subsubsection{2.4. Цилиндрическая
сигма-алгебра}\label{ux446ux438ux43bux438ux43dux434ux440ux438ux447ux435ux441ux43aux430ux44f-ux441ux438ux433ux43cux430-ux430ux43bux433ux435ux431ux440ux430}

На \(E^T\) берут цилиндрическую, или произведенную, \(\sigma\)-алгебру:

\[
\mathcal E^T
=
\sigma\{\pi_{t_1,\ldots,t_n}^{-1}(B):
n\ge1,\ t_i\in T,\ B\in\mathcal B(E^n)\}.
\]

Это минимальная \(\sigma\)-алгебра, относительно которой все
координатные отображения \(X_t\) измеримы.

\subsubsection{2.5. Условия согласованности
распределений}\label{ux443ux441ux43bux43eux432ux438ux44f-ux441ux43eux433ux43bux430ux441ux43eux432ux430ux43dux43dux43eux441ux442ux438-ux440ux430ux441ux43fux440ux435ux434ux435ux43bux435ux43dux438ux439}

Теперь предположим, что для каждого конечного набора попарно различных
индексов

\[
t_1,\ldots,t_n\in T
\]

задана вероятностная мера

\[
\mu_{t_1,\ldots,t_n}
\]

на \(E^n\). Она должна играть роль распределения вектора

\[
(X_{t_1},\ldots,X_{t_n}).
\]

Но произвольный набор мер не подойдет. Нужны условия согласованности.

Первое условие - согласованность по перестановкам. Если \(\sigma\) -
перестановка множества \(\{1,\ldots,n\}\), то распределение координат в
новом порядке должно быть тем же самым распределением, только с
переставленными координатами. Для любого \(B\in\mathcal B(E^n)\):

\[
\mu_{t_{\sigma(1)},\ldots,t_{\sigma(n)}}(B)
=
\mu_{t_1,\ldots,t_n}
\{(x_1,\ldots,x_n):(x_{\sigma(1)},\ldots,x_{\sigma(n)})\in B\}.
\]

Проще: если поменяли порядок индексов, нужно так же поменять порядок
координат.

Второе условие - маргинальная согласованность. Если к набору индексов
добавили лишние моменты времени, а потом не смотрим на соответствующие
координаты, то должна получиться старая мера. Для \(m\ge1\):

\[
\mu_{t_1,\ldots,t_n,t_{n+1},\ldots,t_{n+m}}
(B\times E^m)
=
\mu_{t_1,\ldots,t_n}(B),
\qquad
B\in\mathcal B(E^n).
\]

Особенно важный частный случай:

\[
\mu_{t_1,\ldots,t_n,t_{n+1}}(B\times E)
=
\mu_{t_1,\ldots,t_n}(B).
\]

Это означает, что удаление координаты соответствует взятию маргинального
распределения.

\subsubsection{2.6. Теорема
Колмогорова}\label{ux442ux435ux43eux440ux435ux43cux430-ux43aux43eux43bux43cux43eux433ux43eux440ux43eux432ux430}

Теперь формулировка теоремы.

\textbf{Теорема Колмогорова о согласованных конечномерных
распределениях.} Пусть для всех конечных наборов индексов
\(t_1,\ldots,t_n\in T\) задано семейство вероятностных мер

\[
\mu_{t_1,\ldots,t_n}
\]

на \(E^n\), где \(E\) - стандартное борелевское пространство, например
\(E=\mathbb R\). Если это семейство согласовано по перестановкам и
маргинализации, то существует единственная вероятностная мера \(P\) на
цилиндрической \(\sigma\)-алгебре пространства \(E^T\) такая, что для
всех \(n\), всех \(t_1,\ldots,t_n\) и всех \(B\in\mathcal B(E^n)\):

\[
P\{\pi_{t_1,\ldots,t_n}\in B\}
=
P\{x:(x(t_1),\ldots,x(t_n))\in B\}
=
\mu_{t_1,\ldots,t_n}(B).
\]

\subsubsection{2.7. Канонический координатный
процесс}\label{ux43aux430ux43dux43eux43dux438ux447ux435ux441ux43aux438ux439-ux43aux43eux43eux440ux434ux438ux43dux430ux442ux43dux44bux439-ux43fux440ux43eux446ux435ux441ux441}

Если положить

\[
X_t(x)=x(t),
\]

то семейство

\[
\{X_t:t\in T\}
\]

является случайным процессом на вероятностном пространстве

\[
(E^T,\mathcal E^T,P),
\]

и его конечномерные распределения совпадают с заданными мерами
\(\mu_{t_1,\ldots,t_n}\).

Такой процесс называется каноническим координатным процессом.
``Канонический'' означает, что исход \(\omega\) сам является траекторией
\(x(\cdot)\), а случайная величина \(X_t\) просто берет у этой
траектории значение в момент \(t\).

\subsubsection{2.8. Идея доказательства и роль
согласованности}\label{ux438ux434ux435ux44f-ux434ux43eux43aux430ux437ux430ux442ux435ux43bux44cux441ux442ux432ux430-ux438-ux440ux43eux43bux44c-ux441ux43eux433ux43bux430ux441ux43eux432ux430ux43dux43dux43eux441ux442ux438}

Идея доказательства такая. Сначала хотим задать вероятность на цилиндрах
формулой

\[
P\left(\pi_{t_1,\ldots,t_n}^{-1}(B)\right)
=
\mu_{t_1,\ldots,t_n}(B).
\]

Здесь сразу возникает вопрос корректности: одно и то же цилиндрическое
множество можно записать разными способами. Например,

\[
\pi_t^{-1}(A)
\]

можно также записать как

\[
\pi_{t,s}^{-1}(A\times E).
\]

Чтобы обе записи давали одну и ту же вероятность, нужна маргинальная
согласованность:

\[
\mu_{t,s}(A\times E)=\mu_t(A).
\]

Если поменяли порядок координат, нужна согласованность по перестановкам.
Поэтому условия теоремы ровно гарантируют, что вероятность цилиндра
определена однозначно.

Дальше на классе цилиндров получается предмера. По теореме продолжения
меры, идея которой была в Билете 03, эту предмеру продолжают на
\(\sigma\)-алгебру, порожденную цилиндрами. Полученная мера \(P\) и есть
закон процесса на \(E^T\).

Уникальность тоже естественна: цилиндрические множества порождают
\(\mathcal E^T\), а значения меры на порождающем классе уже заданы
конечномерными распределениями. Поэтому другая мера с теми же
конечномерными распределениями совпадет с \(P\) на всей цилиндрической
\(\sigma\)-алгебре.

\subsubsection{2.9. Ограничения теоремы и свойства
траекторий}\label{ux43eux433ux440ux430ux43dux438ux447ux435ux43dux438ux44f-ux442ux435ux43eux440ux435ux43cux44b-ux438-ux441ux432ux43eux439ux441ux442ux432ux430-ux442ux440ux430ux435ux43aux442ux43eux440ux438ux439}

Важно понимать ограничение теоремы. Она строит процесс на пространстве
всех функций \(E^T\) с цилиндрической \(\sigma\)-алгеброй. Она не
утверждает автоматически, что траектории будут непрерывными,
ограниченными, монотонными или какими-то регулярными. Например, чтобы
получить процесс с непрерывными траекториями, нужны дополнительные
условия регулярности. Теорема Колмогорова дает существование процесса с
заданными конечномерными законами, но не гарантирует хороших свойств
траекторий.

\subsubsection{2.10. Дискретное
время}\label{ux434ux438ux441ux43aux440ux435ux442ux43dux43eux435-ux432ux440ux435ux43cux44f}

В дискретном времени, когда \(T=\mathbb N\), картина особенно проста.
Пространство траекторий:

\[
E^{\mathbb N}.
\]

Цилиндрическое событие задает условия на конечное число координат:

\[
\{x:x_0\in A_0,\ldots,x_n\in A_n\}.
\]

Если все конечномерные распределения согласованы, то существует мера на
пространстве последовательностей, то есть закон случайной
последовательности. Это напрямую ведет к Билету 16.

\subsection{3. Основные
определения}\label{ux43eux441ux43dux43eux432ux43dux44bux435-ux43eux43fux440ux435ux434ux435ux43bux435ux43dux438ux44f}

\subsubsection{Конечномерное распределение
процесса}\label{ux43aux43eux43dux435ux447ux43dux43eux43cux435ux440ux43dux43eux435-ux440ux430ux441ux43fux440ux435ux434ux435ux43bux435ux43dux438ux435-ux43fux440ux43eux446ux435ux441ux441ux430}

Для процесса \(\{X_t:t\in T\}\) и конечного набора \(t_1,\ldots,t_n\)
это распределение случайного вектора

\[
(X_{t_1},\ldots,X_{t_n})
\]

на \(E^n\):

\[
\mu_{t_1,\ldots,t_n}(B)
=
P\{(X_{t_1},\ldots,X_{t_n})\in B\}.
\]

\subsubsection{Пространство
траекторий}\label{ux43fux440ux43eux441ux442ux440ux430ux43dux441ux442ux432ux43e-ux442ux440ux430ux435ux43aux442ux43eux440ux438ux439}

Множество всех функций из \(T\) в \(E\):

\[
E^T=\{x:T\to E\}.
\]

Для дискретной последовательности это \(E^{\mathbb N}\).

\subsubsection{Координатное
отображение}\label{ux43aux43eux43eux440ux434ux438ux43dux430ux442ux43dux43eux435-ux43eux442ux43eux431ux440ux430ux436ux435ux43dux438ux435}

\[
X_t:E^T\to E,
\qquad
X_t(x)=x(t).
\]

\subsubsection{Конечномерная
проекция}\label{ux43aux43eux43dux435ux447ux43dux43eux43cux435ux440ux43dux430ux44f-ux43fux440ux43eux435ux43aux446ux438ux44f}

\[
\pi_{t_1,\ldots,t_n}:E^T\to E^n,
\qquad
\pi_{t_1,\ldots,t_n}(x)=(x(t_1),\ldots,x(t_n)).
\]

\subsubsection{Цилиндрическое
множество}\label{ux446ux438ux43bux438ux43dux434ux440ux438ux447ux435ux441ux43aux43eux435-ux43cux43dux43eux436ux435ux441ux442ux432ux43e}

\[
\pi_{t_1,\ldots,t_n}^{-1}(B),
\qquad
B\in\mathcal B(E^n).
\]

Это событие, зависящее только от конечного числа координат.

\subsubsection{\texorpdfstring{Цилиндрическая
\(\sigma\)-алгебра}{Цилиндрическая \textbackslash sigma-алгебра}}\label{ux446ux438ux43bux438ux43dux434ux440ux438ux447ux435ux441ux43aux430ux44f-sigma-ux430ux43bux433ux435ux431ux440ux430}

\[
\mathcal E^T
=
\sigma\{\pi_{t_1,\ldots,t_n}^{-1}(B):
B\in\mathcal B(E^n)\}.
\]

\subsubsection{Согласованное семейство
мер}\label{ux441ux43eux433ux43bux430ux441ux43eux432ux430ux43dux43dux43eux435-ux441ux435ux43cux435ux439ux441ux442ux432ux43e-ux43cux435ux440}

Семейство конечномерных мер \(\mu_{t_1,\ldots,t_n}\), удовлетворяющее
двум условиям: согласованность по перестановкам и согласованность по
маргинализации.

\subsubsection{Канонический
процесс}\label{ux43aux430ux43dux43eux43dux438ux447ux435ux441ux43aux438ux439-ux43fux440ux43eux446ux435ux441ux441}

Процесс на пространстве \(E^T\), заданный координатами:

\[
X_t(x)=x(t).
\]

\subsection{4. Основные теоремы и
свойства}\label{ux43eux441ux43dux43eux432ux43dux44bux435-ux442ux435ux43eux440ux435ux43cux44b-ux438-ux441ux432ux43eux439ux441ux442ux432ux430}

\subsubsection{Согласованность по
перестановкам}\label{ux441ux43eux433ux43bux430ux441ux43eux432ux430ux43dux43dux43eux441ux442ux44c-ux43fux43e-ux43fux435ux440ux435ux441ux442ux430ux43dux43eux432ux43aux430ux43c}

Для любой перестановки \(\sigma\) координаты можно переставлять только
вместе с индексами:

\[
\mu_{t_{\sigma(1)},\ldots,t_{\sigma(n)}}(B)
=
\mu_{t_1,\ldots,t_n}
\{x:(x_{\sigma(1)},\ldots,x_{\sigma(n)})\in B\}.
\]

\subsubsection{Маргинальная
согласованность}\label{ux43cux430ux440ux433ux438ux43dux430ux43bux44cux43dux430ux44f-ux441ux43eux433ux43bux430ux441ux43eux432ux430ux43dux43dux43eux441ux442ux44c}

Удаление координат соответствует взятию маргинального распределения:

\[
\mu_{t_1,\ldots,t_n,t_{n+1},\ldots,t_{n+m}}
(B\times E^m)
=
\mu_{t_1,\ldots,t_n}(B).
\]

\subsubsection{Теорема
Колмогорова}\label{ux442ux435ux43eux440ux435ux43cux430-ux43aux43eux43bux43cux43eux433ux43eux440ux43eux432ux430-1}

В стандартном борелевском случае всякое согласованное семейство
конечномерных распределений задает единственную вероятностную меру на
цилиндрической \(\sigma\)-алгебре \(E^T\).

\subsubsection{Координатный процесс имеет заданные конечномерные
распределения}\label{ux43aux43eux43eux440ux434ux438ux43dux430ux442ux43dux44bux439-ux43fux440ux43eux446ux435ux441ux441-ux438ux43cux435ux435ux442-ux437ux430ux434ux430ux43dux43dux44bux435-ux43aux43eux43dux435ux447ux43dux43eux43cux435ux440ux43dux44bux435-ux440ux430ux441ux43fux440ux435ux434ux435ux43bux435ux43dux438ux44f}

\[
P\{(X_{t_1},\ldots,X_{t_n})\in B\}
=
\mu_{t_1,\ldots,t_n}(B).
\]

\subsubsection{\texorpdfstring{Единственность закона на цилиндрической
\(\sigma\)-алгебре}{Единственность закона на цилиндрической \textbackslash sigma-алгебре}}\label{ux435ux434ux438ux43dux441ux442ux432ux435ux43dux43dux43eux441ux442ux44c-ux437ux430ux43aux43eux43dux430-ux43dux430-ux446ux438ux43bux438ux43dux434ux440ux438ux447ux435ux441ux43aux43eux439-sigma-ux430ux43bux433ux435ux431ux440ux435}

Если две меры на \(E^T\) совпадают на всех цилиндрах, то они совпадают
на всей \(\sigma\)-алгебре, порожденной цилиндрами.

\subsubsection{Теорема не дает регулярность
траекторий}\label{ux442ux435ux43eux440ux435ux43cux430-ux43dux435-ux434ux430ux435ux442-ux440ux435ux433ux443ux43bux44fux440ux43dux43eux441ux442ux44c-ux442ux440ux430ux435ux43aux442ux43eux440ux438ux439}

Непрерывность, правосторонняя непрерывность, ограниченность и другие
свойства траекторий требуют отдельных условий.

\subsection{5. Примеры}\label{ux43fux440ux438ux43cux435ux440ux44b}

\subsubsection{Независимая одинаково распределенная
последовательность}\label{ux43dux435ux437ux430ux432ux438ux441ux438ux43cux430ux44f-ux43eux434ux438ux43dux430ux43aux43eux432ux43e-ux440ux430ux441ux43fux440ux435ux434ux435ux43bux435ux43dux43dux430ux44f-ux43fux43eux441ux43bux435ux434ux43eux432ux430ux442ux435ux43bux44cux43dux43eux441ux442ux44c}

Пусть \(\nu\) - вероятностная мера на \(E\). Для конечного набора разных
индексов \(t_1,\ldots,t_n\in\mathbb N\) зададим

\[
\mu_{t_1,\ldots,t_n}
=
\nu^{\otimes n}.
\]

То есть все координаты имеют одно и то же распределение \(\nu\) и
независимы. Это семейство согласовано: перестановка координат не меняет
произведение одинаковых мер, а маргинализация произведения дает
произведение меньшего числа множителей:

\[
\nu^{\otimes(n+1)}(B\times E)=\nu^{\otimes n}(B).
\]

По теореме Колмогорова существует случайная последовательность
независимых одинаково распределенных величин с законом \(\nu\).

\subsubsection{Последовательность независимых, но не одинаково
распределенных
величин}\label{ux43fux43eux441ux43bux435ux434ux43eux432ux430ux442ux435ux43bux44cux43dux43eux441ux442ux44c-ux43dux435ux437ux430ux432ux438ux441ux438ux43cux44bux445-ux43dux43e-ux43dux435-ux43eux434ux438ux43dux430ux43aux43eux432ux43e-ux440ux430ux441ux43fux440ux435ux434ux435ux43bux435ux43dux43dux44bux445-ux432ux435ux43bux438ux447ux438ux43d}

Пусть для каждого \(k\in\mathbb N\) задана мера \(\nu_k\) на \(E\). Для
индексов \(k_1,\ldots,k_n\) задаем

\[
\mu_{k_1,\ldots,k_n}
=
\nu_{k_1}\otimes\cdots\otimes\nu_{k_n}.
\]

Это семейство тоже согласовано. Значит можно построить
последовательность независимых величин с заранее заданными разными
одномерными распределениями.

\subsubsection{Марковская цепь по начальному распределению и переходной
матрице}\label{ux43cux430ux440ux43aux43eux432ux441ux43aux430ux44f-ux446ux435ux43fux44c-ux43fux43e-ux43dux430ux447ux430ux43bux44cux43dux43eux43cux443-ux440ux430ux441ux43fux440ux435ux434ux435ux43bux435ux43dux438ux44e-ux438-ux43fux435ux440ux435ux445ux43eux434ux43dux43eux439-ux43cux430ux442ux440ux438ux446ux435}

Пусть пространство состояний счетное, начальное распределение равно
\(\pi\), а переходная матрица равна \(P=(p_{ij})\). Для
\(i_0,\ldots,i_n\) задаем конечномерные вероятности:

\[
\mu_{0,\ldots,n}(i_0,\ldots,i_n)
=
\pi_{i_0}p_{i_0i_1}p_{i_1i_2}\cdots p_{i_{n-1}i_n}.
\]

Эти распределения согласованы: если просуммировать по последней
координате, то из условия

\[
\sum_j p_{ij}=1
\]

получается распределение первых \(n\) координат. Поэтому теорема
Колмогорова дает существование марковской цепи с заданными \(\pi\) и
\(P\). Это связь с Билетом 22 и Билетом 23.

\subsubsection{Гауссовский
процесс}\label{ux433ux430ux443ux441ux441ux43eux432ux441ux43aux438ux439-ux43fux440ux43eux446ux435ux441ux441}

Чтобы задать гауссовский процесс, задают среднюю функцию

\[
m(t)
\]

и ковариационную функцию

\[
K(s,t).
\]

Для каждого конечного набора \(t_1,\ldots,t_n\) берут гауссовское
распределение в \(\mathbb R^n\) со средним вектором

\[
(m(t_1),\ldots,m(t_n))
\]

и ковариационной матрицей

\[
(K(t_i,t_j))_{i,j=1}^n.
\]

Чтобы это было корректно, все такие матрицы должны быть неотрицательно
определены. Тогда гауссовские конечномерные распределения согласованы и
по теореме Колмогорова существует гауссовский процесс. Это связь с
Билетом 18.

\subsubsection{Пример
несогласованности}\label{ux43fux440ux438ux43cux435ux440-ux43dux435ux441ux43eux433ux43bux430ux441ux43eux432ux430ux43dux43dux43eux441ux442ux438}

Пусть пытаются задать

\[
P\{X_1=0\}=1
\]

и одновременно

\[
P\{(X_1,X_2)=(1,1)\}=1.
\]

Такого процесса не существует, потому что маргинал двумерного
распределения по первой координате дает

\[
P\{X_1=1\}=1,
\]

что противоречит заданному одномерному распределению. Это показывает,
зачем нужна согласованность.

\subsection{6. Схемы доказательств /
обоснований}\label{ux441ux445ux435ux43cux44b-ux434ux43eux43aux430ux437ux430ux442ux435ux43bux44cux441ux442ux432-ux43eux431ux43eux441ux43dux43eux432ux430ux43dux438ux439}

\subsubsection{Почему условия согласованности
необходимы}\label{ux43fux43eux447ux435ux43cux443-ux443ux441ux43bux43eux432ux438ux44f-ux441ux43eux433ux43bux430ux441ux43eux432ux430ux43dux43dux43eux441ux442ux438-ux43dux435ux43eux431ux445ux43eux434ux438ux43cux44b}

Если процесс уже существует, то для любого конечного набора

\[
(X_{t_1},\ldots,X_{t_n})
\]

его распределение не должно зависеть от того, как мы записали одно и то
же событие.

Если поменяли порядок координат, событие не изменилось по смыслу,
поменялась только запись. Значит нужна согласованность по перестановкам.

Если из вектора

\[
(X_{t_1},\ldots,X_{t_n},X_s)
\]

не смотреть на последнюю координату, то должно получиться распределение

\[
(X_{t_1},\ldots,X_{t_n}).
\]

Значит нужна маргинальная согласованность.

\subsubsection{Идея построения
меры}\label{ux438ux434ux435ux44f-ux43fux43eux441ux442ux440ux43eux435ux43dux438ux44f-ux43cux435ux440ux44b}

\begin{enumerate}
\def\labelenumi{\arabic{enumi}.}
\tightlist
\item
  На цилиндрах задаем:
\end{enumerate}

\[
P(\pi_{t_1,\ldots,t_n}^{-1}(B))
=
\mu_{t_1,\ldots,t_n}(B).
\]

\begin{enumerate}
\def\labelenumi{\arabic{enumi}.}
\setcounter{enumi}{1}
\item
  Согласованность гарантирует, что это определение не зависит от
  представления цилиндра.
\item
  Получаем предмеру на алгебре цилиндров.
\item
  По теореме продолжения меры продолжаем ее на \(\sigma\)-алгебру,
  порожденную цилиндрами.
\item
  На полученном вероятностном пространстве берем координатный процесс:
\end{enumerate}

\[
X_t(x)=x(t).
\]

\subsubsection{Почему закон
единственен}\label{ux43fux43eux447ux435ux43cux443-ux437ux430ux43aux43eux43d-ux435ux434ux438ux43dux441ux442ux432ux435ux43dux435ux43d}

Если две вероятностные меры на \(E^T\) имеют одинаковые конечномерные
распределения, то они совпадают на всех цилиндрах. Цилиндры порождают
\(\mathcal E^T\), значит меры совпадают на всей \(\mathcal E^T\).

\subsection{7. Что написать на
доске}\label{ux447ux442ux43e-ux43dux430ux43fux438ux441ux430ux442ux44c-ux43dux430-ux434ux43eux441ux43aux435}

\begin{enumerate}
\def\labelenumi{\arabic{enumi}.}
\tightlist
\item
  Пространство траекторий:
\end{enumerate}

\[
E^T=\{x:T\to E\}.
\]

\begin{enumerate}
\def\labelenumi{\arabic{enumi}.}
\setcounter{enumi}{1}
\tightlist
\item
  Проекция:
\end{enumerate}

\[
\pi_{t_1,\ldots,t_n}(x)=(x(t_1),\ldots,x(t_n)).
\]

\begin{enumerate}
\def\labelenumi{\arabic{enumi}.}
\setcounter{enumi}{2}
\tightlist
\item
  Цилиндры:
\end{enumerate}

\[
\pi_{t_1,\ldots,t_n}^{-1}(B),
\qquad
B\in\mathcal B(E^n).
\]

\begin{enumerate}
\def\labelenumi{\arabic{enumi}.}
\setcounter{enumi}{3}
\tightlist
\item
  Согласованность по маргинализации:
\end{enumerate}

\[
\mu_{t_1,\ldots,t_n,s}(B\times E)
=
\mu_{t_1,\ldots,t_n}(B).
\]

\begin{enumerate}
\def\labelenumi{\arabic{enumi}.}
\setcounter{enumi}{4}
\tightlist
\item
  Согласованность по перестановкам:
\end{enumerate}

\[
\mu_{t_{\sigma(1)},\ldots,t_{\sigma(n)}}
=
\mu_{t_1,\ldots,t_n}\circ p_\sigma^{-1},
\]

где

\[
p_\sigma(x_1,\ldots,x_n)=(x_{\sigma(1)},\ldots,x_{\sigma(n)}).
\]

\begin{enumerate}
\def\labelenumi{\arabic{enumi}.}
\setcounter{enumi}{5}
\tightlist
\item
  Вывод теоремы:
\end{enumerate}

\[
\exists!P\ \text{на}\ (E^T,\mathcal E^T):
\quad
P\{\pi_{t_1,\ldots,t_n}\in B\}
=
\mu_{t_1,\ldots,t_n}(B).
\]

Здесь \(\mathcal E^T\) - цилиндрическая \(\sigma\)-алгебра.

\begin{enumerate}
\def\labelenumi{\arabic{enumi}.}
\setcounter{enumi}{6}
\tightlist
\item
  Канонический процесс:
\end{enumerate}

\[
X_t(x)=x(t).
\]

\subsection{8. Типичные дополнительные
вопросы}\label{ux442ux438ux43fux438ux447ux43dux44bux435-ux434ux43eux43fux43eux43bux43dux438ux442ux435ux43bux44cux43dux44bux435-ux432ux43eux43fux440ux43eux441ux44b}

\textbf{Что является главным объектом теоремы?}

Семейство конечномерных распределений

\[
\{\mu_{t_1,\ldots,t_n}\}.
\]

Теорема говорит, когда оно является семейством конечномерных
распределений некоторого процесса.

\textbf{Почему одних одномерных распределений недостаточно?}

Потому что они не задают зависимость между разными моментами времени.
Например, две величины могут иметь одинаковые одномерные распределения,
но быть независимыми, полностью зависимыми или иметь любую допустимую
совместную структуру. Нужны все совместные распределения конечных
наборов координат.

\textbf{Что такое цилиндрическое событие?}

Это событие, зависящее только от конечного числа координат траектории:

\[
\{x:(x(t_1),\ldots,x(t_n))\in B\}.
\]

Остальные координаты не ограничиваются.

\textbf{Зачем нужна согласованность по перестановкам?}

Чтобы распределение не зависело от порядка записи одного и того же
конечного набора координат. Вектор \((X_s,X_t)\) и вектор \((X_t,X_s)\)
имеют меры, связанные перестановкой координат.

\textbf{Зачем нужна маргинальная согласованность?}

Чтобы распределение меньшего набора координат совпадало с маргиналом
распределения большего набора. Например:

\[
P\{(X_s,X_t)\in B\times E\}
=
P\{X_s\in B\}.
\]

\textbf{Что именно строит теорема Колмогорова?}

Она строит вероятностную меру \(P\) на пространстве траекторий \(E^T\) с
цилиндрической \(\sigma\)-алгеброй. После этого координаты
\(X_t(x)=x(t)\) становятся случайным процессом.

\textbf{Гарантирует ли теорема непрерывные траектории?}

Нет. Она гарантирует только существование процесса с заданными
конечномерными распределениями. Регулярность траекторий требует
дополнительных условий.

\textbf{В чем разница между законом процесса и конечномерными
распределениями?}

Закон процесса - это мера на всем пространстве траекторий \(E^T\).
Конечномерные распределения - это ее маргиналы при конечномерных
проекциях.

\textbf{Где здесь используется теорема продолжения меры?}

Сначала мера задается на цилиндрах, то есть на событиях конечномерного
вида. Затем ее продолжают на \(\sigma\)-алгебру, порожденную цилиндрами.

\textbf{Как связана теорема с Билетом 16?}

В Билете 16 дискретный процесс рассматривается как случайная
последовательность, то есть как случайный элемент пространства
\(E^{\mathbb N}\). Теорема Колмогорова как раз дает меру на этом
пространстве по согласованным конечномерным распределениям.

\subsection{9. Типичные
ошибки}\label{ux442ux438ux43fux438ux447ux43dux44bux435-ux43eux448ux438ux431ux43aux438}

\begin{itemize}
\tightlist
\item
  Задавать только одномерные распределения и считать, что процесс уже
  определен.
\item
  Забывать маргинальную согласованность.
\item
  Забывать согласованность по перестановкам.
\item
  Путать конечномерное распределение и траекторию.
\item
  Путать закон процесса на \(E^T\) и отдельные распределения на \(E^n\).
\item
  Считать, что теорема автоматически дает непрерывные или регулярные
  траектории.
\item
  Не указывать, на какой \(\sigma\)-алгебре строится мера.
\item
  Формулировать теорему без условий на пространство состояний; для
  аккуратного ответа безопасно говорить о \(E=\mathbb R\) или
  стандартном борелевском пространстве.
\end{itemize}

\subsection{10. Связи с другими
билетами}\label{ux441ux432ux44fux437ux438-ux441-ux434ux440ux443ux433ux438ux43cux438-ux431ux438ux43bux435ux442ux430ux43cux438}

\begin{itemize}
\tightlist
\item
  Билет 03: идея продолжения меры с алгебры цилиндров на порожденную
  \(\sigma\)-алгебру.
\item
  Билет 05: координатный процесс состоит из измеримых отображений.
\item
  Билет 06: конечномерные распределения процесса - распределения
  случайных векторов.
\item
  Билет 16: дискретный процесс как случайный элемент пространства
  последовательностей \(E^{\mathbb N}\).
\item
  Билет 17: стационарность в узком смысле формулируется через
  конечномерные распределения.
\item
  Билет 18: гауссовский процесс задается согласованными гауссовскими
  конечномерными распределениями.
\item
  Билет 22 и Билет 23: марковские цепи строятся по начальному
  распределению и переходным вероятностям.
\end{itemize}

\subsection{11. Минимум на
удовлетворительно}\label{ux43cux438ux43dux438ux43cux443ux43c-ux43dux430-ux443ux434ux43eux432ux43bux435ux442ux432ux43eux440ux438ux442ux435ux43bux44cux43dux43e}

Нужно сказать:

\begin{enumerate}
\def\labelenumi{\arabic{enumi}.}
\tightlist
\item
  Конечномерное распределение процесса - это закон вектора
\end{enumerate}

\[
(X_{t_1},\ldots,X_{t_n}).
\]

\begin{enumerate}
\def\labelenumi{\arabic{enumi}.}
\setcounter{enumi}{1}
\item
  Чтобы набор таких распределений задавал процесс, он должен быть
  согласован.
\item
  Два условия согласованности:
\end{enumerate}

\[
\text{перестановка индексов} \Longleftrightarrow \text{перестановка координат},
\]

и

\[
\mu_{t_1,\ldots,t_n,s}(B\times E)
=
\mu_{t_1,\ldots,t_n}(B).
\]

\begin{enumerate}
\def\labelenumi{\arabic{enumi}.}
\setcounter{enumi}{3}
\item
  Тогда существует процесс на \(E^T\) с этими конечномерными
  распределениями.
\item
  Главная ловушка: теорема не дает регулярность траекторий.
\end{enumerate}

\subsection{12. Минимум на
хорошо}\label{ux43cux438ux43dux438ux43cux443ux43c-ux43dux430-ux445ux43eux440ux43eux448ux43e}

Помимо минимума, нужно объяснить:

\begin{itemize}
\tightlist
\item
  что такое пространство траекторий \(E^T\);
\item
  что такое цилиндрическое событие;
\item
  что такое координатный процесс \(X_t(x)=x(t)\);
\item
  почему согласованность нужна для корректного задания меры на
  цилиндрах;
\item
  как идея доказательства связана с продолжением меры;
\item
  привести пример iid-последовательности или марковской цепи.
\end{itemize}

\subsection{13. Что добавить для
отлично}\label{ux447ux442ux43e-ux434ux43eux431ux430ux432ux438ux442ux44c-ux434ux43bux44f-ux43eux442ux43bux438ux447ux43dux43e}

Для сильного ответа стоит добавить:

\begin{itemize}
\tightlist
\item
  аккуратную формулировку через конечномерные проекции
  \(\pi_{t_1,\ldots,t_n}\);
\item
  единственность меры на цилиндрической \(\sigma\)-алгебре;
\item
  различие между законом процесса и конечномерными распределениями;
\item
  предупреждение про отсутствие автоматической регулярности траекторий;
\item
  пример несогласованного семейства;
\item
  пример гауссовского процесса через среднюю и ковариационную функции;
\item
  связь с узкой стационарностью, марковскими цепями и случайными
  последовательностями.
\end{itemize}

\subsection{14. Устная версия ответа на 5
минут}\label{ux443ux441ux442ux43dux430ux44f-ux432ux435ux440ux441ux438ux44f-ux43eux442ux432ux435ux442ux430-ux43dux430-5-ux43cux438ux43dux443ux442}

\begin{enumerate}
\def\labelenumi{\arabic{enumi}.}
\tightlist
\item
  Теорема Колмогорова нужна, чтобы по конечномерным распределениям
  построить случайный процесс.
\item
  Пусть для каждого конечного набора \(t_1,\ldots,t_n\) задана мера
  \(\mu_{t_1,\ldots,t_n}\) на \(E^n\).
\item
  Эти меры должны быть согласованы по перестановкам: если меняем порядок
  индексов, меняется порядок координат.
\item
  Они должны быть согласованы по маргинализации:
\end{enumerate}

\[
\mu_{t_1,\ldots,t_n,s}(B\times E)
=
\mu_{t_1,\ldots,t_n}(B).
\]

\begin{enumerate}
\def\labelenumi{\arabic{enumi}.}
\setcounter{enumi}{4}
\tightlist
\item
  Рассматриваем пространство траекторий
\end{enumerate}

\[
E^T=\{x:T\to E\}.
\]

\begin{enumerate}
\def\labelenumi{\arabic{enumi}.}
\setcounter{enumi}{5}
\tightlist
\item
  На нем есть координатные отображения
\end{enumerate}

\[
X_t(x)=x(t).
\]

\begin{enumerate}
\def\labelenumi{\arabic{enumi}.}
\setcounter{enumi}{6}
\tightlist
\item
  Цилиндры имеют вид
\end{enumerate}

\[
\{x:(x(t_1),\ldots,x(t_n))\in B\}.
\]

\begin{enumerate}
\def\labelenumi{\arabic{enumi}.}
\setcounter{enumi}{7}
\tightlist
\item
  По заданным конечномерным мерам задаем вероятность цилиндра:
\end{enumerate}

\[
P(\pi_{t_1,\ldots,t_n}^{-1}(B))
=
\mu_{t_1,\ldots,t_n}(B).
\]

\begin{enumerate}
\def\labelenumi{\arabic{enumi}.}
\setcounter{enumi}{8}
\tightlist
\item
  Согласованность делает это определение корректным.
\item
  Потом мера продолжается на \(\sigma\)-алгебру, порожденную цилиндрами.
\item
  Полученный координатный процесс имеет ровно заданные конечномерные
  распределения.
\item
  Важно: теорема дает существование закона процесса, но не гарантирует
  непрерывность или другие свойства траекторий.
\end{enumerate}

\subsection{15. Если осталось 2 минуты перед
экзаменом}\label{ux435ux441ux43bux438-ux43eux441ux442ux430ux43bux43eux441ux44c-2-ux43cux438ux43dux443ux442ux44b-ux43fux435ux440ux435ux434-ux44dux43aux437ux430ux43cux435ux43dux43eux43c}

Запомнить:

\[
E^T=\{x:T\to E\},
\qquad
X_t(x)=x(t).
\]

Цилиндр:

\[
\pi_{t_1,\ldots,t_n}^{-1}(B).
\]

Главное условие:

\[
\mu_{t_1,\ldots,t_n,s}(B\times E)
=
\mu_{t_1,\ldots,t_n}(B),
\]

плюс согласованность по перестановкам.

Главная формулировка: согласованное семейство конечномерных
распределений задает единственную меру на пространстве траекторий с
цилиндрической \(\sigma\)-алгеброй, а координаты становятся процессом.

Главная ловушка: теорема не утверждает регулярность траекторий; она
только строит закон процесса.

\newpage

\end{document}
